\pdfoutput=1
\documentclass[11pt]{article}
\usepackage[margin=1in]{geometry}
\usepackage{amsmath,amssymb,amsthm}
\usepackage{enumitem}
\usepackage{booktabs}
\usepackage{url}
\usepackage{hyperref}
\hypersetup{hidelinks,
  pdftitle={An initial Theta-datum over an imaginary quadratic field, and numerical values of the local quantities of inter-universal Teichmuller theory IV},
  pdfauthor={Toshihisa Urahigashi}}

\theoremstyle{plain}
\newtheorem{theorem}{Theorem}
\newtheorem{proposition}[theorem]{Proposition}
\newtheorem{corollary}[theorem]{Corollary}
\theoremstyle{definition}
\newtheorem{hypothesis}{Hypothesis}
\newtheorem{remark}{Remark}
\newtheorem{question}{Question}

\newcommand{\Fmod}{F_{\mathrm{mod}}}
\newcommand{\Vmod}{\mathbb{V}_{\mathrm{mod}}}
\newcommand{\Vbad}{\underline{\mathbb{V}}^{\mathrm{bad}}}
\newcommand{\Vgood}{\mathbb{V}^{\mathrm{good}}}
\newcommand{\Vgoodmod}{\mathbb{V}^{\mathrm{good}}_{\mathrm{mod}}}
\newcommand{\VV}{\mathbb{V}}
\newcommand{\Ind}[1]{\textup{(Ind#1)}}
\newcommand{\SA}{\textup{(SA)}}
\newcommand{\ord}{\operatorname{ord}}
\newcommand{\Norm}{\operatorname{N}}
\newcommand{\qbb}{\underline{\underline q}}
\newcommand{\VVu}{\underline{\mathbb V}}

\title{An initial $\Theta$-datum over an imaginary quadratic field,\\
and numerical values of the local quantities of\\
inter-universal Teichm\"uller theory IV\\
\large A second report}
\author{Toshihisa Urahigashi\\
\small\texttt{123026.urahigashi.toshihisa@gmail.com}\\[2pt]
\small Version 1.0.4\\[1pt]
\small DOI \href{https://doi.org/10.5281/zenodo.22601333}{10.5281/zenodo.22601333} (all versions);\\[1pt]
\small this version \href{https://doi.org/10.5281/zenodo.22634210}{10.5281/zenodo.22634210}}
\date{8 September 2026\\[3pt]
\small ORCID \href{https://orcid.org/0009-0004-0460-6242}{\texttt{0009-0004-0460-6242}}}

\begin{document}
\maketitle

\begin{abstract}
We exhibit an initial $\Theta$-datum in the sense of \cite[Def.~3.1]{IUT1} whose
field of moduli is the \emph{imaginary quadratic} field $\mathbb{Q}(i)$: the curve
$E_0:y^2=x(x-1)(x-a)$ with $a=(10+3i)^{29}$ and $\ell=7$, for which the rational
prime $109$ splits, so that one place of multiplicative and one place of good
reduction lie above it.  All conditions (a)--(f) are verified; the arithmetic is
exact, the geometric conditions (d)--(f) are established in \S\ref{s:geom}, and
the computations, including the separate exhaustive power-freeness sieve, are
described in \S\ref{s:repro}.  The local invariants are $e_{\mathrm g}=1$,
$e_{\mathrm b}=105$ and
$\ord_{109}(\underline{\underline q}_{\mathrm b})=29/7$, so that the
\emph{ramification} invariants of this datum agree with those of the datum over
$\mathbb{Q}(\sqrt5)$ in \cite{First}.  The residual data do not: the residue degrees
are $f_{\mathrm g}=12$, $f_{\mathrm b}=6$ here against $24$ and $2$ there, and the mixed
tensor algebra has $6$ components of degree $1260$ against $2$ of degree $2520$
\emph{(the values at $109$ are computed in \path{scripts/tables.py}, Table 13, from
the point count and the order of Frobenius on $E_0[3]$, $E_0[5]$, $E_0[7]$; the values
at $211$ are those of \cite{First}.  The consistency check
$12\cdot630=7560=6\cdot1260$ is asserted there.)}  Everything tabulated in
\S\ref{s:orders} is computed at $211$ and does not transfer.

Explicit initial $\Theta$-data are needed for any numerically effective use of the
theory.  Dupuy and Hilado \cite{DH25} construct such data systematically for elliptic
curves over $\mathbb{Q}$, working out the curve of Cremona label 11a1 in full and
reporting $5525$ curves of prime conductor for which the mod-$\ell$ representation
over $\mathbb Q(i,E[30])$ is surjective for some $\ell\in\{7,11,13\}$
\cite[\S2.11, p.~16]{DH25}.  This count does not certify all the conditions of
initial $\Theta$-data for all those curves.  The datum below is over
$\mathbb{Q}(i)$; the point of it is not that an explicit datum is rare but that the
\emph{configuration} it carries is unavailable over $\mathbb{Q}$.  It is also, in three
specific respects, easier to construct than the datum of \cite{First}: the nodes are
already split over the field of moduli, the degree bound is halved, and the existence
family needs two congruences instead of five, with the condition on $\ell$ becoming
automatic (\S\ref{s:easier}).  We give that family as
Proposition~\ref{p:family}: for every $N$ there are infinitely many \emph{arithmetic
parameters} of this shape over $\mathbb{Q}(i)$ with $\ord_\pi(a)=N$, so the
configuration is not isolated.  \emph{Turning one of them into an initial
$\Theta$-datum needs more}: $7\nmid N$ (else $\ell=7$ divides the Tate order $2N$), the
Galois image condition, and the geometric choices of \S\ref{s:geom}; the proposition is
an arithmetic statement and does not assert them.

The effective results of \cite{MFHMP} are stated for a mono-complex field, that is,
for $\mathbb{Q}$ or an imaginary quadratic field \cite[Thm.~A]{MFHMP}, with an
explicit constant in the degree-two case, and the curve there is in the same Legendre
form; the datum of \cite{First} lies over a real quadratic field, and the one given
here does not.  \emph{Whether that bears on the estimate is not addressed here.}

Explicit data also make the quantities in the local estimates of
\cite[Thm.~1.10]{IUT4} numerical at these inputs.  \emph{Explicit formulas, and the
plan of evaluating them, are already in \cite{IUT4}, \cite{MFHMP} and \cite{DHPS}};
what is added here is their exact evaluation on the two data and a set of checks that
detects the kind of error \S\ref{s:regime} records.  Sections
\ref{s:margin}--\ref{s:slack} tabulate them for this datum and for that of
\cite{First}.  Every entry is either a formula printed in the source, marked
\textup{[S]}, or a computation from it, marked \textup{[P]}; each table names
the locator of the quantity it tabulates, and every number is reproduced by
\path{scripts/tables.py}, which uses only the standard library --- with one exception,
named where it occurs: the counts of the Cremona sweep in \S\ref{s:regime} are
transcribed from an external computation, and $\ord_3(q)=44$ at the two curves it
returns is taken from that computation as an input, the arithmetic that follows from
$44$ being what is recomputed.  Three of the tables
evaluate, on these explicit inputs, three questions: where the data on record sit
relative to the threshold at which the stage-one comparison ceases to be vacuous
(\S\ref{s:regime}); which region of the parameters $(\ell,p,N)$ admits a datum
satisfying at once that threshold, the choice of $\ell$ made in
\cite[Cor.~2.2(ii)]{IUT4}, and the condition $I^{*}=\varnothing$ of
\cite[Prop.~1.4(iii)]{IUT4} (\S\ref{s:region}); and how the two terms of the averaged
upper bound compare with each other at these parameters (\S\ref{s:slack}).  For a
candidate that already satisfies the other three, that region is a window
$\sqrt{H_\forall}\le\ell<N/3$, its left wall the lower half of the source's condition
(C1) and its right wall the threshold derived here --- the \emph{upper} half of (C1) is
checked separately and is not binding on any datum tabulated here --- and
\S\ref{s:fourteen} exhibits a single explicit curve over $\mathbb{Q}(i)$ whose window
contains \emph{fourteen} primes, with two distinct upper constraints
simultaneously tight at $\ell=157$.
\emph{Scope}: this report constructs data and evaluates local quantities on them.  It
does not compute the hull of the union of the possible images of the actual
$\Theta$-pilot or compare the two sides of \cite[Cor.~3.12]{IUT3}; nothing here is
conditional on the hypothesis \SA{} of \cite{First}; and no claim is made about
\cite[Cor.~3.12]{IUT3}, about \cite[Thm.~A]{MFHMP}, about the objections of
\cite{SS}, or about $abc$.
\end{abstract}

\medskip
\noindent\textbf{2020 Mathematics Subject Classification.}
Primary 14H25; Secondary 11G05, 14G32, 11R58.

\noindent\textbf{Keywords.}
inter-universal Teichm\"uller theory, initial $\Theta$-datum, imaginary quadratic
field, mono-complex field, Legendre curve, split prime, mixed reduction.

\section{What this is, and what it is for}\label{s:why}

\textup{[P]} An initial $\Theta$-datum is a list of conditions
\cite[Def.~3.1]{IUT1} that a curve, a prime $\ell$, a set of bad places, a section
and a cusp must jointly satisfy.  Verifying the list against an explicit curve is a
finite task, and doing it produces an object that others can compute with.
\textup{[S]} \cite[\S2.10]{DH25} does this over $\mathbb{Q}$: it breaks
\cite[Def.~3.1]{IUT1} into pieces, works out the curve 11a1 in full, and reports that
for $5525$ curves of prime conductor in the Cremona tables some
$\ell\in\{7,11,13\}$ makes the mod-$\ell$ representation over $\mathbb{Q}(i,E[30])$
surjective \cite[\S2.11, p.~16]{DH25} --- that one condition, not the whole
of \cite[Def.~3.1]{IUT1}.
Zhou also describes constructions of $\mu_6$-initial $\Theta$-data over
$\mathbb Q$, including Frey--Hellegouarch curves
\cite[Prop.~2.5, Lemma~2.6]{Zhou}.
These are rational-moduli constructions.
\textup{[P]} \emph{Over $\mathbb{Q}$ the mixed
profile cannot appear at all}: by \cite[Cor.~3(i)]{First} every place tuple above every
prime is constant when $\Fmod=\mathbb{Q}$, so no member of that family carries it,
and the field must be at least quadratic for the configuration to exist.  \cite{First}
gives a datum over $\mathbb{Q}(\sqrt5)$ that does carry it; the present report gives
one over $\mathbb{Q}(i)$, and shows that the construction is in several respects
easier there (\S\ref{s:easier}).  The value of this does not
depend on any reading of the theory, on any conclusion about $abc$, or on the outcome
of any dispute: a verified datum is a verified datum.

\textup{[P]} Two further things come out of the construction.  First, the
configuration in which a rational prime splits in the field of moduli and carries one
place of multiplicative and one of good reduction occurs here, as it does in
\cite{First}; by \cite[Cor.~3(i)]{First} it cannot occur when
$\Fmod=\mathbb{Q}$, so the field must be at least quadratic for it to appear at all.
Second, it is not isolated: Proposition~\ref{p:family} produces infinitely many
arithmetic parameters of this shape over $\mathbb{Q}(i)$, of arbitrary depth at the
split prime --- subject, for the passage to a datum, to $7\nmid N$ and to the conditions
of \S\ref{s:geom}.

\begin{remark}[where such data are used]\label{r:mfhmp}
\textup{[S]} \cite[Thm.~A]{MFHMP} is stated for ``a mono-complex number field
[i.e., $\mathbb{Q}$ or an imaginary quadratic field --- cf.\ Definition 1.2]'',
Definition 1.2 there reading ``$F$ is mono-complex if $F$ is either the field of
rational numbers $\mathbb{Q}$ or an imaginary quadratic field''.  It provides the
constant
\[
 h_d(\epsilon)=
 \begin{cases}
  3.4\cdot10^{30}\,\epsilon^{-166/81} & (d=1),\\
  6\cdot10^{31}\,\epsilon^{-174/85} & (d=2),
 \end{cases}
 \qquad d=[L:\mathbb{Q}],
\]
so the degree-two case is provided for explicitly; and that the curve there is
$E_{a,b,c}:y^2=x(x-1)(x+a/c)$, the Legendre form used here.  The datum of
\cite{First} lies over a real quadratic field and so is not of that kind; the datum
below is.  \emph{That is a statement about where the two data lie, and nothing more:
whether the configuration bears on the estimate of \cite[Thm.~A]{MFHMP} is an open
question which this report does not address and on which none of its claims depend.}
\end{remark}

\begin{remark}[independence]\label{r:indep}
\textup{[P]} \cite{SS}, \cite{Rpt} and \cite{LANA} are listed because a reader will
want to know where this sits relative to them; \emph{no statement below uses any of
them, and no claim below is conditional on how the questions they are about are
answered}.  In particular nothing here depends on the hypothesis \SA{} of
\cite{First}.  What is not computed is the hull of the union of the possible images
of the actual $\Theta$-pilot, and what is not compared is the two sides of
\cite[Cor.~3.12]{IUT3}: what is asserted is a construction and its verification.
\end{remark}

\section{The datum}\label{s:datum}

\begin{equation}\label{e:datum}
\begin{aligned}
 &L=\Fmod=\mathbb{Q}(i), &&\pi=10+3i,\quad \Norm(\pi)=109,\quad a=\pi^{29},\\
 &E_0:\,y^2=x(x-1)(x-a), &&\ell=7,\\
 &F=L(E_0[15]), &&K=F(E_0[7])=L(E_0[105]).
\end{aligned}
\end{equation}
Let $X_F=E_F\setminus\{0\}$.  Select as $\Vbad_{\mathrm{mod}}$ all places of $L$ of
multiplicative reduction outside residue characteristics $2$ and $7$, including the
place $v_\pi$.

\begin{theorem}\label{t:datum}
\textup{[P]} \eqref{e:datum} extends to an initial $\Theta$-datum in the sense of
\cite[Def.~3.1]{IUT1}.  Above $109$ the set $\VV$ has exactly two elements, one
selected bad and one good, and
\[
 e_{\mathrm g}=1,\qquad e_{\mathrm b}=105,\qquad
 \ord_{109}(\underline{\underline q}_{\mathrm b})=\frac{29}{7},
\]
so every local index in this profile satisfies $e\le p-2=107$.
\end{theorem}

The arithmetic conditions are verified in \S\ref{s:arith} and the geometric ones in
\S\ref{s:geom}.  \emph{The slack $107-105=2$ is small, but the inequality $105\le107$ is what is
required; \S\ref{s:variants} lists roomier primes.}

\section{The arithmetic conditions}\label{s:arith}

The finite integer, rational and residue-field computations of this section are exact
and are reproduced by \path{scripts/imq_checks.py}.  \emph{Two things in this section
are not}: the seventh-power-free sieve, which is a separate program over all primes
below its bound (\S\ref{s:repro}), and the appeals to cited structure theorems, which
are applications and not computations.

\paragraph{(a) The field of moduli.}
Writing $\bar a$ for the conjugate of $a$, the six Legendre orbit comparisons
$\bar a\in\{a,a^{-1},1-a,(1-a)^{-1},a(a-1)^{-1},(a-1)a^{-1}\}$ all fail in
$\mathbb{Z}[i]$.  Hence $j(E_0)\ne\overline{j(E_0)}$ and
$\Fmod=\mathbb{Q}(j(E_0))=\mathbb{Q}(i)$.  Computing $j$ exactly in
$\mathbb{Q}(i)$, its imaginary part, in lowest terms, is a nonzero
integer of $296$ digits over a denominator of $234$ digits.

\paragraph{The mixed profile above $109$.}
Since $109=\pi\bar\pi$ splits and $a=\pi^{29}$,
\[
 \ord_\pi(a)=29,\qquad \ord_{\bar\pi}(a)=0,\qquad
 \ord_\pi(1-a)=\ord_{\bar\pi}(1-a)=0 .
\]
So $E_0$ has multiplicative reduction at $\pi$ and good reduction at $\bar\pi$, and
$\VV$ has one selected bad and one good element above $109$.

\paragraph{(b) Selected places, and the extra hypothesis of \cite{IUT4}.}
$\Norm(a)=109^{29}$ and
\[
 \Norm(1-a)=2\cdot3^2\cdot5\cdot21577\cdot C,
\]
with $C$ of $53$ digits.  The full norm is seventh-power-free: the exhaustive sieve
of \path{scripts/imq_sieve7.cpp} tests every prime up to
$\lfloor C^{1/7}\rfloor=34{,}869{,}670$ and finds none.  Primality of $C$ is not
asserted.  Every selected place has odd residue characteristic, and
$7\nmid\Norm(1-a)$.  Since the bad set consists of \emph{all} multiplicative
places outside $2$ and $7$, the additional good-reduction hypothesis of
\cite[Thm.~1.10]{IUT4} holds by construction.

\paragraph{(b) Stable reduction.}
\textup{[S]} \cite[Def.~3.1(b)]{IUT1} asks for stable reduction at every
$v\in\mathbb V(F)^{\mathrm{non}}$.  \cite[Prop.~1.8(v)]{IUT4} gives it for a curve all
of whose $l$-torsion is rational, for one prime $l\ge3$ \emph{invertible in
$\mathcal{O}_k$}.  \textup{[P]} Here $E_0[15]\subset E_0(F)$ by construction, so both
$l=3$ and $l=5$ are available; a nonarchimedean place has a single residue
characteristic, so at least one of them is invertible there.  \emph{That is what the
choice of $15=3\cdot5$ buys}: neither $3$ nor $5$ alone would cover its own residue
characteristic.  Marking the origin and contracting the components that this
destabilises carries the semistable model of $E_0$ to a stable model of the marked
curve of type $(1,1)$, which is what \cite[Def.~3.1(b)]{IUT1} asks for.  No further
condition on $a$ is needed.

\paragraph{(c) The order at every selected place, and the mod-$\ell$ image.}
By the norm identity
$v_q(\Norm(z))=\sum_{\mathfrak p\mid q}f(\mathfrak p/q)v_{\mathfrak p}(z)$,
seventh-power-freeness gives $v_{\mathfrak p}(1-a)\le6$ at every prime ideal, hence
prime to $7$.  At $v_\pi$ the base Tate order is $2\cdot29=58$ and $7\nmid58$.

For the image: at $v_\pi$ the Tate inertia contains the transvection
$T=\bigl(\begin{smallmatrix}1&1\\0&1\end{smallmatrix}\bigr)$, nontrivial because
$7\nmid58$.  Each rational prime $q\equiv1\pmod4$ has \emph{two} places of
$\mathbb{Q}(i)$ above it, and they \emph{may} give different Frobenius traces, so a
witness is a place and not merely a prime.  (They need not: at $q=233$ both places
send $a$ to $12$, so both give trace $-26$.)  Take the place at which $i$ reduces to the smaller
root $r$ of $r^2\equiv-1$:
\[
 \begin{array}{r|rrrrrr}
  q          & 13 & 17 & 37 & 41 & 53 & 61\\\hline
  r          &  5 &  4 &  6 &  9 & 23 & 11\\
  a_q        &  6 & -6 & -2 & -6 & -6 & -10\\
  a_q^2-4q\bmod7 &  5 &  3 &  3 &  5 &  6 &  3
 \end{array}
\]
Each discriminant is a nonsquare modulo $7$, so $X^2-a_qX+q$ is irreducible there.
\emph{The choice of place is not cosmetic}: at $q=13$ the other place has $a_q=2$
and discriminant $1$, and at $q=61$ it has $a_q=-6$ and discriminant $2$, both
squares, so those two places are not witnesses.  \textup{[P]} \emph{The passage from these facts to the image is made without choosing
a basis.}  Write $L$ for the fixed line of $T$ and $g$ for a Frobenius at one of the six
places.  Since the characteristic polynomial of $g$ is irreducible, $g$ has no
eigenvector, so $gL\ne L$; hence $T$ and $gTg^{-1}$ are nontrivial transvections with
\emph{distinct} fixed lines.  As $\mathrm{SL}_2(\mathbb F_7)$ is transitive on ordered
pairs of distinct lines, the pair may be moved to $(\langle e_1\rangle,\langle
e_2\rangle)$, where the two transvections read
$\bigl(\begin{smallmatrix}1&a\\0&1\end{smallmatrix}\bigr)$ and
$\bigl(\begin{smallmatrix}1&0\\b&1\end{smallmatrix}\bigr)$ with $a,b\ne0$; conjugating
further by $\mathrm{diag}(a^{-1},1)$ makes $a=1$.  Each of the six remaining pairs
generates a group of order $336=|\mathrm{SL}_2(\mathbb F_7)|$.  Every step is checked in
the kernel in \path{lean/BasisFreeImage.lean}: that all $882$ matrices with irreducible
characteristic polynomial move all eight lines, that the orbit of
$(\langle e_1\rangle,\langle e_2\rangle)$ has all $56$ elements, and the six closures.

\begin{remark}[correction]\label{r:basis}
\textup{[P]} An earlier version conjugated $T$ by the \emph{companion matrix} of the
characteristic polynomial and checked that one pair.  That does not settle the image:
the actual Frobenius is only conjugate to the companion matrix, and the basis putting
inertia in the form $T$ need not be the one putting Frobenius in companion form.  The
argument above quantifies over all matrices with irreducible characteristic polynomial
and over all pairs of distinct lines, so no basis is chosen anywhere.
\end{remark}
  Finally $[F:\Fmod]$ divides
\[
 |\mathrm{GL}_2(\mathbb Z/15\mathbb Z)|=23040=2^9\cdot3^2\cdot5 ,
\]
which is prime to $7$, while $|\mathrm{PSL}_2(\mathbb F_7)|=168=2^3\cdot3\cdot7$ is
divisible by $7$; every proper normal subgroup of $\mathrm{SL}_2(\mathbb F_7)$ lies
in its centre.  So intersecting with the image over $F$ cannot give a proper
subgroup, and $\mathrm{Gal}(K/F)$ has the required image.

\paragraph{The local invariants.}
At $\bar\pi$ the reduction is good and the torsion is prime to $109$, so
$e_{\mathrm g}=1$.  At $\pi$ the curve is split Tate with $\ord_{109}(q)=58$; since
$\gcd(105,58)=1$, adjoining $E_0[105]$ has ramification index exactly $105$, whence
$e_{\mathrm b}=105$, and $\ord_{109}(\underline{\underline q}_{\mathrm b})=58/14=29/7$.

\section{The geometric conditions}\label{s:geom}

\paragraph{Split reduction at every selected bad place.}
$a$ is integral over $\mathbb{Z}[i]$.  At an odd place where $a$ and $a-1$ are both
units the discriminant $16a^2(a-1)^2$ is a unit, so only the two nodal reductions
need inspection.  If $a\equiv0$ the special fibre is $y^2=x^2(x-1)$ with tangent
cone $y^2=-x^2$, which splits because $i\in L$.  If $a\equiv1$, putting $t=x-1$
gives $y^2=(1+t)t^2$ with tangent cone $y^2=t^2$, already split.  The residue
characteristics are odd, so the branches are distinct.  \emph{Reduction is
therefore split multiplicative already over $L$ at every selected bad place, and
stays split over $F$ and $K$.}  This is not deduced from rational $2$-torsion
alone.

\paragraph{(d) Quotient and core.}
Write $V=E_0[7](K)$, fix a line $W\subset V$ and $0\ne\xi\in Q=V/W$, and let
$\pi_W:E_K\to E'=E_K/W$ with $\psi:E'\to E_K$, $\psi\pi_W=[7]$.  Restriction gives
$\bar\pi_W:Q\simeq\ker\psi$, so this kernel is constant over $K$.  Put
$\underline X_K=E'\setminus\ker\psi$,
$\underline C_K=[\underline X_K/\{\pm1\}]$ and $C_K=[X_K/\{\pm1\}]$.  The cover
$\underline X_K\to X_K$ has geometric quotient $Q$; its extension to the
compactifications is unramified above the zero cusp with $K$-rational fibre there,
so the cusp decomposition group maps trivially to $Q$.  Passing to the inversion
quotients gives type $(1,7\text{-tors})^\pm$ as in \cite[Def.~2.1]{EtTh}.  Let
$\epsilon$ be the cusp corresponding to $\pm\bar\pi_W(\xi)$.

By \cite[Prop.~2.1]{MFHMP} an arithmetic once-punctured elliptic curve in
characteristic zero has one of the four \emph{rational} invariants
$2^{14}31^35^{-3}$, $2^2 73^3 3^{-4}$, $1728$, $0$.  The present $j$ is not
rational, so $X_{\overline K}$ --- where $X_K:=E_K\setminus\{0\}$ is the
once-punctured curve --- is nonarithmetic; the finite \'etale map
$X_{\overline K}\to C_{\overline K}$ puts these orbicurves in the same
commensurability class, so $C_{\overline K}$ is nonarithmetic as well.  It is a
punctured hemi-elliptic orbicurve, so \cite[Prop.~2.7]{CanLift} makes it a
$\overline K$-core, and \cite[Prop.~2.3(i)]{CanLift} descends this to $C_K$, which
is then the core of the finite \'etale cover $\underline C_K$ required by
\cite[Def.~3.1(d)]{IUT1}.  For a completion $K_v$, extend $K\hookrightarrow
\overline K$ to $\overline K\hookrightarrow\overline{K_v}$ and apply
\cite[Prop.~2.3(ii)]{CanLift} to these algebraically closed fields of
characteristic zero, then (i) again; base change preserves finite \'etaleness.
This covers the complex completions, and no real completion occurs since $i\in K$.
The geometric theta covers of (d) are supplied by
\cite[Prop.~2.2, Def.~2.3]{EtTh}.  \emph{We note that $j$ is not an algebraic
integer, so $E_0$ has no complex multiplication; but no non-CM criterion is used
above.}

\paragraph{(e) and (f): places and cusp, simultaneously.}
At each selected bad moduli place $u$ choose a place of $F$ above $u$ and a place
$w$ of $K$ above that.  The split Tate sequence
$0\to\mu_7\to E(K_w)[7]\to\mathbb Z/7\to0$, with the class of $q^{1/7}$ mapping to
$1$, determines a marked flag $(W_w,\pm\xi_w)$ in $V$: changing the root moves the
lift only by $\mu_7$, and reversing the graph orientation changes the sign.  The
group $\mathrm{SL}_2(V)$ acts transitively on such flags --- lift a nonzero
quotient class to $t$ and choose $s\in W$ with $\det(s,t)=1$ --- so, by the image
statement of \S\ref{s:arith}, there is $\sigma_u\in\mathrm{Gal}(K/F)$ with
$\sigma_u(W_w,\pm\xi_w)=(W,\pm\xi)$.  Put $v(u)=\sigma_u w$, with the convention
$|b|_{\sigma_u w}=|\sigma_u^{-1}b|_w$; the resulting isomorphism
$\sigma_u^{-1}:K_{v(u)}\simeq K_w$ commutes with localization, so the fixed global
pair acquires the required local marking at $v(u)$.  \emph{The global $W$ and
$\epsilon$ are not replaced place by place.}

Since $\sigma_u$ fixes $F$, the place $v(u)$ lies above the same $u$; distinct
moduli-place fibres are disjoint, so \emph{the choices at different $u$ do not
interfere and can be made simultaneously}.  Choosing one arbitrary lift at each
remaining moduli place gives $\VV\xrightarrow{\ \sim\ }\Vmod$ as required by
\cite[Def.~3.1(e)]{IUT1}, which asks for a set-theoretic section and imposes
neither equivariance nor a single common Galois element.

Locally the quotient is $E(q)\to E(q^7)$, $z\mapsto z^7$, with dual
$E(q^7)\to E(q)$, $z\mapsto z$; the latter multiplies the period of the graph by
$7$.  \emph{Two things have to hold, and we separate them so that each can be checked
against \cite[Def.~2.5]{EtTh}.}  \emph{First}, the fundamental-group quotient factors
through $\widehat{\mathbb Z}\to\mathbb Z/7$, which is the displayed local picture.
\emph{Second}, the chosen cusp splitting is compatible with the $\pm$ structure: since
$\ker\psi=\{q^r\bmod q^{7\mathbb Z}:r\in\mathbb Z/7\}$ and $\pi_W(q^{1/7})=q$, the
localization of the fixed cusp $\epsilon$ \emph{is} the canonical $\pm1$ cusp of
\cite[Def.~3.1(f)]{IUT1}: the label counts deck translations by the period $q$,
and no Tate order or ramification index enters it.  Rational full $2$-torsion gives
the field condition of \cite[Def.~3.1(e)]{IUT1}; the Weil pairings give
$\mu_{12}\subset F$ and $\mu_7\subset K$.

\section{An arithmetic existence family}\label{s:family}

The configuration is not isolated.  Put
$f(T)=(T^2+1)((T-1)^2+1)=T^4-2T^3+3T^2-2T+2$, of discriminant
$400=2^4 5^2$.

\begin{proposition}\label{p:family}
\textup{[P]} Let $p$ be a rational prime that splits in $\mathbb{Q}(i)$ with
$p\ne5$ --- equivalently here $p\ge13$ --- and let $N$ be a positive integer.  There are infinitely many positive integers $A$ such that,
for $a=A-i$ and the two places $\pi,\bar\pi$ of $\mathbb{Q}(i)$ above $p$,
\[
 \ord_\pi(a)=N,\qquad \ord_{\bar\pi}(a)=0,\qquad
 \ord_\pi(1-a)=\ord_{\bar\pi}(1-a)=0,
\]
and $\Norm(a)\Norm(1-a)/p^N$ is a positive seventh-power-free integer.
\end{proposition}

\begin{proof}[Construction]
Hensel gives $r_N$ with $r_N^2\equiv-1\pmod{p^{N+1}}$ and $r_N\equiv r_1\pmod p$.
Impose
\[
 A\equiv r_N+p^{N}\pmod{p^{N+1}},\qquad A\equiv0\ \text{or}\ 1\pmod5,
\]
put $M=5p^{N+1}$, $A=A_0+Mx$ and $G_N(x)=p^{-N}f(A_0+Mx)$, a quartic in $x$.  \emph{The
exclusion of $p=5$ is needed for the two congruences to be simultaneously solvable}: if
$p=5$ then $r_N\equiv\pm2\pmod5$, so $A\equiv r_N+5^{N}\equiv\pm2\pmod5$ for $N\ge1$,
which meets neither $A\equiv0$ nor $A\equiv1$; the moduli $5$ and $p^{N+1}$ must also be
coprime for the two to be combined.  The
first congruence forces $\ord_\pi(a)=N$ exactly; since $\pi\mid(A-i)$ implies
$\pi\nmid(A+i)$, this is $v_p(A^2+1)$.  \emph{The prime $p$ itself does not divide the
other factor}: $A\equiv r_1\pmod p$, so
\[
 (A-1)^2+1\equiv(r_1-1)^2+1=1-2r_1\pmod p ,
\]
and if this vanished then $2r_1\equiv1$, whence $4r_1^2\equiv1$ and, using
$r_1^2\equiv-1$, $p\mid5$ --- excluded.  Hence $v_p(f(A))=v_p(A^2+1)=N$ and
$p\nmid G_N(x)$.

Only two congruences are needed, because of the following three facts, each checked
exhaustively.  \emph{(i)} $v_2(f(A))=1$ for every $A$: if $A$ is odd then
$A^2+1\equiv2\pmod 8$ and $(A-1)^2+1$ is odd, and if $A$ is even the roles are
exchanged.  So no condition modulo $2$ is required and $2^7\nmid f(A)$ always.
\emph{(ii)} If $q\equiv3\pmod4$ then $A^2+1\equiv0\pmod q$ has no solution, so
$q\nmid f(A)$; this removes $3,7,11,19,23,\dots$ at once.  \emph{(iii)} In particular $7\nmid f(A)$ for every $A$, so
\emph{no selected bad place has residue characteristic $7$}.  \textup{[P]}
\emph{This is not the same as the order condition of \cite[Def.~3.1(c)]{IUT1}}, which
asks that $\ell$ not divide $v_{\mathfrak p}(a)$; that is secured at the places away
from $\pi$ by seventh-power-freeness, and at $\pi$ itself only by the extra requirement
$7\nmid N$, since the Tate order there is $2N$.  \emph{These are orders at places of
$L$, and \cite[Def.~3.1(c)]{IUT1} asks about places of $F$}: for $w\mid v$ the Tate
order is multiplied by $e(w/v)$, and since $F=L(E_0[15])$ is Galois over $L$ with
$[F:L]\mid|\mathrm{GL}_2(\mathbb Z/15)|=23040=2^{9}3^{2}5$, every $e(w/v)$ is prime to
$7$.  So an order prime to $7$ downstairs stays prime to $7$ at every bad place of $F$.  The second congruence makes
$v_5(f(A))=0$.

Hence a prime $q$ with $q^{7}\mid G_N(x)$ satisfies $q\equiv1\pmod4$ and $q\ge13$.
\emph{This is a statement about seventh powers, not about divisors}: $2$ divides
$G_N(x)$, since $v_2(f(A))=1$ always.  Such a $q$ is neither $5$ nor $p$, by the two
paragraphs above, and is not $2$; \emph{therefore $q\nmid M=5p^{N+1}$}, so multiplication
by $M$ is invertible modulo $q^{7}$ and the roots in $x$ correspond bijectively to the
roots of $f$.  Such $q$
are prime to $\operatorname{disc}(f)=400$, so $f$ has at most four simple roots
modulo $q$, each lifting uniquely modulo $q^7$.  With
$0<G_N(x)<C_N(X+1)^4$ for $1\le x\le X$ and $C_N=M^4/p^N$, the number of excluded
$x$ is at most
\[
 \sum_{\substack{q\ \mathrm{prime},\ 13\le q\le Q_N\\ q\equiv1\ (4)}}
 4\Bigl(\frac{X}{q^7}+1\Bigr)
 \le 7.38\times10^{-8}\,X+4C_N^{1/7}(X+1)^{4/7},
 \qquad Q_N:=\bigl\lfloor(C_N(X+1)^4)^{1/7}\bigr\rfloor,
\]
whose second term is $o(X)$ for fixed $N$ and whose first coefficient is far below
$1$.  Infinitely many $x$ therefore remain.  \emph{The constant is proved, not
estimated}: summing $4q^{-7}$ exactly over the primes $q\equiv1\pmod4$ with
$13\le q\le100$ and bounding the rest by $4\int_{100}^{\infty}x^{-7}\,dx=2/(3\cdot10^{12})$
--- which over-counts, since it runs over all integers --- gives
$7.3794978287\ldots\times10^{-8}<7.38\times10^{-8}$ in exact rational arithmetic
(\path{scripts/imq_family.py}).  \emph{A truncated positive sum is only a lower bound},
so the earlier floating-point check did not establish this inequality.

\textup{[P]} \emph{The range matters}: written over all $q$ the term $+1$ would sum
to infinity.  Only $q\le Q_N$ can produce an exclusion, so the $+1$ contributions are
counted over that finite range; the main terms $X/q^{7}$ are then extended to the full
sum, which only increases them.  The density bound is unaffected.
\end{proof}

\emph{This is an arithmetic existence statement.  It does not verify (d)--(f) for
each member; those are the geometric conditions of \S\ref{s:geom}, established for
the datum \eqref{e:datum}.}

\section{\texorpdfstring{Three ways in which $\mathbb{Q}(i)$ is simpler}{Three ways in which Q(i) is simpler}}\label{s:easier}

\textup{[P]} Compared with the datum of \cite{First} over $\mathbb{Q}(\sqrt5)$:

\begin{enumerate}[label=(\roman*),itemsep=2pt]
\item \emph{The node is already split.}  Over $\mathbb{Q}_{211}$ the element $-1$ is
 a nonsquare, so \cite{First} adjoins $i$ to split the bad node.  Here $109\equiv1
 \pmod4$ and $i\in\Fmod$, so both nodal types split over $L$ itself
 (\S\ref{s:geom}).
\item \emph{The degree bound is halved.}  Write
 $G=|\mathrm{GL}_2(\mathbb Z/15\mathbb Z)|=23040$.  Here $[F:\Fmod]$ divides $G$,
 whereas in \cite{First} it divides $2G=46080$: no factor of $2$ is spent on
 adjoining $i$.
\item \emph{The family needs two congruences instead of five.}  The five conditions
 of the family in \cite{First} reduce here to the two of
 Proposition~\ref{p:family}, and $\operatorname{disc}(f)=400$ rather than
 $144400=20^2 19^2$.  Comparing the two excluded densities \emph{as the same
 quantity} --- the sum $\sum4/q^7$ over the primes each family must exclude, not the
 integral majorant $1/69984$ that \cite{First} uses in its counting step --- the
 value falls from $4.93\times10^{-6}$ to $7.38\times10^{-8}$, a factor of $66.8$.
 One term accounts for almost all of it: every $q\equiv3\pmod4$ is excluded here
 automatically, \emph{including $q=\ell=7$}, and $4/7^7=4.86\times10^{-6}$ is
 $98.5\%$ of the density for \cite{First}.  So the structural gain and the numerical
 gain are the same fact: the condition on $\ell$ becomes automatic rather than a
 separate check.
\end{enumerate}

\section{Roomier variants}\label{s:variants}

\textup{[P]} The constraint that keeps $I^{*}$ empty is $e_{\mathrm b}\le p-2$ with
\[
 e_{\mathrm b}=\frac{15\ell}{\gcd(15\ell,2N)},
\]
$p$ split in $\Fmod$.  \emph{It reduces to $e_{\mathrm b}=15\ell$ only when the gcd is
$1$}, as it is at $(N,\ell)=(29,7)$, where $\gcd(105,58)=1$; the figures below are for
that fixed pair.  Varying $N$ changes the gcd --- at $(114,37)$ it is $3$ --- and also
requires $\ell\nmid2N$ to be rechecked.  For $\ell=7$ and $N=29$ this asks
$p\ge107$; the datum above uses the first such prime, so the slack $p-2-e_{\mathrm b}$
is $2$.  Larger split primes give more: $113$, $137$, $149$, $157$, $173$, $181$,
$193$ leave $6,30,42,50,66,74,86$.  Every imaginary quadratic field of class number
one admits a split prime $\ge107$.  We have kept $p=109$, and $N=29$, so that the
local invariants agree exactly with \cite{First}.  \emph{$N$ is not free}: moving it
changes $\gcd(15\ell,2N)$, hence $e_{\mathrm b}$, and the condition $\ell\nmid2N$ of
\cite[Def.~3.1(c)]{IUT1} has to hold as well.

\section{The stage margins and their thresholds}\label{s:margin}

\textup{[P]} With $I^{*}=\varnothing$ admissible in \cite[Prop.~1.4(iii)]{IUT4}, and
\emph{for a tuple whose last entry is bad} --- the case in which $\lambda=j^2s$ rather
than $0$, by Step (v) of the proof of \cite[Thm.~1.10]{IUT4} --- the difference
$\mathrm{co}-\mathrm{bd}$ at stage $j$ is $j^2s-(j+2)$, and it is positive exactly when
$\ord_p(q)>2\ell(j+2)/j^2$.  \emph{These are evaluations of the coefficients in the
bound}, not a comparison of the actual $\Theta$-pilot with anything.

\begin{center}
\begin{tabular}{cccc}
\toprule
$j$ & $j^2s$ & $j^2s-(j+2)$ & positive iff $\ord_p(q)>$\\
\midrule
1 & $29/7$  & $8/7$   & $42$\\
2 & $116/7$ & $88/7$  & $14$\\
3 & $261/7$ & $226/7$ & $70/9$\\
\bottomrule
\end{tabular}
\end{center}

\textup{[P]} \emph{The binding stage is $j=1$, by a factor of three over $j=2$ and of
more than five over $j=3$.}  Everything about the threshold is therefore decided at
the first stage.

\section{Where the data on record sit}\label{s:regime}

\textup{[P]} $\ord_p(q)=-v_p(j_{E})$ at a place of multiplicative reduction.  At
$\ell=7$ the stage-one threshold of \S\ref{s:margin} reads $\ord_p(q)>42$.

\begin{center}
\begin{tabular}{lrrrr}
\toprule
datum & $p$ & $\ord_p(q)$ & $s$ & $s-3$\\
\midrule
\cite{DH25}, Cremona 11a1 over $\mathbb{Q}$ & 11 & 5 & $5/14$ & $-2.64$\\
Cremona 14a1 over $\mathbb{Q}$ & 7 & 3 & $3/14$ & $-2.79$\\
Cremona 15a1 over $\mathbb{Q}$ & 5 & 4 & $2/7$  & $-2.71$\\
Cremona 37a1 over $\mathbb{Q}$ & 37 & 1 & $1/14$ & $-2.93$\\
\cite{First} over $\mathbb{Q}(\sqrt5)$ & 211 & \textbf{58} & $29/7$ & $\mathbf{+1.14}$\\
the present datum over $\mathbb{Q}(i)$ & 109 & \textbf{58} & $29/7$ & $\mathbf{+1.14}$\\
\bottomrule
\end{tabular}
\end{center}

\textup{[P]} \emph{The last column of the four $\mathbb{Q}$-rows is the value of the
formula, not a licensed comparison}: at 14a1 $p=\ell$
(which \cite[Def.~3.1(c)]{IUT1} does not allow at a bad place at all), and at 15a1
$p\mid15$.  At 11a1 and 37a1 the exclusion is by size, and \emph{the bound does not
depend on choosing a field}: \cite[Def.~3.1(b)]{IUT1} requires the $2\cdot3$-torsion
of $E_F$ to be rational over $F$, so at the reference level $\ell=7$ the field $K$
contains $E[42]$.  Both primes are prime to $42$, and their Tate orders are $5$ and
$1$, so
\[
 e_{\mathrm b}\ \ge\ \frac{42}{\gcd(42,\ord_p(q))}=42
 \quad\text{at both},\qquad 42>9,\qquad 42>35 .
\]
\emph{An earlier version printed $105/\gcd(105,5)=21$ at 11a1 and $105$ at 37a1},
carrying over the torsion level $105$ of the datum \eqref{e:datum} to curves that do
not have it; \cite{DH25} takes $F=\mathbb{Q}(i,E[30])$, which at level $7$ would give
$210/\gcd(210,5)=42$ and at its own level $13$ gives $390/\gcd(390,5)=78$.  The
exclusions are unchanged; the printed indices were not the indices of these curves.
They are printed to show the scale of $\ord_p(q)$ on record, nothing more.

\textup{[P]} The sample above is not the whole record.  Sweeping the Cremona tables
of conductor below $10^4$ --- $64687$ curves --- for a \emph{single} place of
multiplicative reduction with $\ord_p(q)>42$ returns $80$ curves.  Of these, $78$ have
that place above $2$, which \cite[Def.~3.1]{IUT1} excludes: bad places of residue
characteristic two are not admitted.  The remaining two, Cremona \texttt{5160j1} and
\texttt{5190c3}, have the place above $3$ with $\ord_3(q)=44$, so that at $\ell=7$

\[
  s=\tfrac{44}{14}=\tfrac{22}{7},\qquad s-3=\tfrac17>0 .
\]

\textup{[P]} So the threshold \emph{is} reached inside the tabulated range, and it is
reached by curves that have been on record since long before this line of work.  What
those two curves do not do is admit the simplified form of the estimate: with $p-2=1$
and $e\ge2$ (see below), those places must be put into the set $I^{*}$ of
\cite[Prop.~1.4(iii)]{IUT4}.  \emph{The estimate still applies}: Step (v)
of \cite[Thm.~1.10]{IUT4} is written for exactly this case and absorbs such places into
$4(j+1)\,\ell^{*}_{\mathrm{mod}}$.  What is not available at those curves is the
$I^{*}=\varnothing$ form, which is the one the margin of \S\ref{s:margin} is computed
from; the margin is therefore not the applicable quantity there.

\textup{[P]} \emph{No number is put on $e$ at those two curves.}  The prime-to-$p$
expression $15\ell/\gcd(15\ell,\ord_p(q))$ is not available when $p$
divides the torsion level, and here $p=3$ divides $105$.  What can be said without it is
enough: full $3$-torsion forces $\mu_3\subset K$ by the Weil pairing, and
$\zeta_3-1$ is a root of the Eisenstein polynomial $T^2+3T+3$, so
$\mathbb{Q}_3(\mu_3)/\mathbb{Q}_3$ is totally ramified of index $2$ and $e\ge2>1=p-2$.
\emph{This also settles the number that the unavailable expression would have
returned}: $2\mid e$, so $e=105$ is not merely unproved at those places but impossible.  The threshold of \S\ref{s:margin} is therefore not by itself
a useful filter; it is one of five conditions, and \S\ref{s:region} treats them
together.  The two data of \cite{First} and above were built with $a=\pi^{29}$, which
forces $\ord_p(q)=58$ at a prime large enough for $e_{\mathrm b}\le p-2$ to hold as
well.

\section{Tuple types}\label{s:tuple}

\textup{[S]} \cite[Thm.~1.10, Step (v)]{IUT4} fixes $\lambda$ by the status of the
\emph{last} index alone, verbatim: ``$\lambda$'' to be $0$ if $\underline
v_j\in\underline{\mathbb V}^{\mathrm{good}}$; ``$\lambda$'' to be ``$\ord(-)$'' of the
element $\qbb_{\underline v_j}^{\,j^2}$ [\dots] if $\underline v_j\in\Vbad$.
\emph{The $q$ carries the double underline}: it is the $2\ell$-th root of
\cite[Ex.~3.2(iv)]{IUT1}, not $q$ itself, which is why $\lambda=j^2s$ with
$s=\ord_p(\qbb)=29/7$ and not $j^2\ord_p(q)=58j^2$.  \textup{[P]} Write $d_I=(j+1)-\sum_i e_{t_i}^{-1}$, and put
\[
  \mathrm{co}(t):=-\sum_i e_{t_i}^{-1},\qquad
  \mathrm{bd}(t):=-\lambda+d_I+1 .
\]
\emph{$\mathrm{co}(t)$ is the coefficient expression, not the log-volume of the
tensor packet}, and the two differ once a tuple has more than one bad entry.  The good
side is unramified, so a factor $R_{\mathrm g}$ costs nothing.  With $k$ bad entries
$R_I$ is non-maximal \emph{exactly when $k\ge2$}; for $k=0,1$ it is maximal.  The packet
$\log_p(R_I^{\times})$ therefore has normalized log-volume
$\mathrm{co}(t)\log p+\mu^{\log}(R_I)$, that is, \emph{below} $\mathrm{co}(t)\log p$ by
\[
  -\mu^{\log}(R_I)=\tfrac{\max\{k-1,0\}(e-1)}{2e}\log p,\qquad e=e_{\mathrm b}=105,
\]
which is $\tfrac{52}{105}\log p$ for $k=2$ --- the gain recorded in \S\ref{s:orders} ---
and $\tfrac{104}{105}\log p$ for $k=3$.  So the rows unaffected are those with
\emph{at most one} bad entry; among the rows printed below, $(\mathrm b,\mathrm b)$,
$(\mathrm g,\mathrm b,\mathrm b)$ and $(\mathrm b,\mathrm b,\mathrm g)$ carry the
$k=2$ defect and $(\mathrm b,\mathrm b,\mathrm b)$ the $k=3$ one.  \emph{Nothing in the
table depends on this}: the columns are $\mathrm{co}$ and $\mathrm{bd}$ as defined.  (Earlier drafts called $\mathrm{co}$ the ``seed'',
which collided with the $\lceil\lambda\rceil$ of \S\ref{s:container}; the two are
different quantities and are now named apart.)

\begin{center}
\begin{tabular}{clrrr}
\toprule
$j$ & tuple & $\mathrm{co}(t)$ & $\mathrm{bd}(t)$ & $\mathrm{co}-\mathrm{bd}$\\
\midrule
1 & $(\mathrm g,\mathrm g)$ & $-2$ & $1$ & $-3$\\
1 & $(\mathrm g,\mathrm b)$ & $-106/105$ & $-226/105$ & $\mathbf{8/7}$\\
1 & $(\mathrm b,\mathrm g)$ & $-106/105$ & $209/105$ & $-3$\\
1 & $(\mathrm b,\mathrm b)$ & $-2/105$ & $-122/105$ & $\mathbf{8/7}$\\
\midrule
2 & $(\mathrm g,\mathrm g,\mathrm g)$ & $-3$ & $1$ & $-4$\\
2 & $(\mathrm g,\mathrm b,\mathrm b)$ & $-107/105$ & $-1427/105$ & $\mathbf{88/7}$\\
2 & $(\mathrm b,\mathrm b,\mathrm g)$ & $-107/105$ & $313/105$ & $-4$\\
2 & $(\mathrm b,\mathrm b,\mathrm b)$ & $-1/35$ & $-63/5$ & $\mathbf{88/7}$\\
\bottomrule
\end{tabular}
\end{center}

\textup{[P]} \emph{The difference is constant on each of two classes, not on all
types.}  It is $8/7$ at $j=1$ and $88/7$ at $j=2$ exactly for the tuples whose last
entry is bad, and $-3$, $-4$ for those whose last entry is good.  In particular
``mixed'' does not name a value: $(\mathrm g,\mathrm b)$ and $(\mathrm b,\mathrm g)$
have the same $\mathrm{co}$ and different bounds.  The cancellation behind the constant value on
the first class is $d_I+\sum_ie_{t_i}^{-1}=j+1$; see \cite[\S6]{First}.

\begin{remark}[correction]\label{r:lambda}
\textup{[P]} An earlier version of this table set $\lambda=j^{2}s$ for \emph{every}
type, including the all-good ones, and reported a single ``mixed'' row.  Both are
corrected above.  The all-good rows change from $(-22/7,\,8/7)$ and $(-109/7,\,88/7)$ to
$(1,-3)$ and $(1,-4)$; the rows for tuples ending in a bad place are unchanged.
\end{remark}

\section{How much probability each tuple carries}\label{s:prob}

\textup{[S]} \cite{DHPS} puts a probability measure on the tuples, and the tables
above are per-tuple quantities, so it is worth recording what weight each of them
receives on the present data.  The measure is stated twice; we quote both as printed.  In \cite[\S1,
p.~2]{DHPS} the numerator is the local degree of the \emph{base} field,
$\Pr(v)=[F_{0,v}:\mathbb{Q}_p]/[F_0:\mathbb{Q}_p]$; in \cite[\S7.1,
p.~22]{DHPS} the denominator is $[F_0:\mathbb{Q}]^{j+1}$.  Read together the two
statements fix the measure, and the degree formula
$\sum_{v\mid p}[F_{0,v}:\mathbb{Q}_p]=[F_0:\mathbb{Q}]$ then makes it a probability
measure:
\[
  \Pr(v_0,\dots,v_j)=\frac{2}{\ell-1}\prod_{i=0}^{j}
  \frac{[F_{0,v_i}:\mathbb{Q}_p]}{[F_0:\mathbb{Q}]}.
\]

\textup{[P]} For both data of \S\ref{s:datum} the prime splits in $F_0$, so there
are two places above it, each of local degree $1$, and $[F_0:\mathbb{Q}]=2$.  With
$\ell=7$ the index $j$ runs over $1,2,3$:

\begin{center}
\begin{tabular}{rrrr}
\toprule
$j$ & tuples $2^{j+1}$ & $\Pr$ per tuple & mass at this $j$\\
\midrule
1 & 4  & $1/12$ & $1/3$\\
2 & 8  & $1/24$ & $1/3$\\
3 & 16 & $1/48$ & $1/3$\\
\midrule
 & 28 & & $1$\\
\bottomrule
\end{tabular}
\end{center}

\textup{[P]} So the mixed $j=1$ tuple of \S\ref{s:tuple} --- the one carrying the
$8/7$ of that table --- carries $1/12$ of the mass, and each index $j$ contributes
$1/3$ regardless of how many tuples it has.  Two remarks on reading the source
here.  \emph{We adopt a reading for the sole purpose of evaluating the displayed
formula, and make no claim about what either point does to the arguments of that
paper.}  The
numerator in \cite[\S7.1]{DHPS} is printed with a completion degree of $K$ rather
than of $F_0$; taken literally the tuples would carry masses well above $1$, and it
is the \S1 statement that identifies the field.  We do not consider what either point
does to the estimates of that paper; we record the reading only because the tables
here evaluate that very formula on explicit data, so the normalisation has to be
pinned down, and we use it as the degree normalisation of
\cite[Rem.~1.7.1]{IUT4}.

\section{The region cut out by five conditions}\label{s:region}

Five conditions can be asked of a datum in this family.  \emph{Only (2) and (5) are
hypotheses the source imposes.}  (1) is the threshold derived in \S\ref{s:margin};
(3) is a non-vacuity criterion of ours; and (4), as explained below, is not a
hypothesis at all but the condition under which a place may be left out of the
set $I^{*}$ of \cite[Prop.~1.4(iii)]{IUT4}.

\begin{enumerate}[label=(\arabic*),itemsep=1pt]
\item $\ord_p(q)=2N>6\ell$ \hfill \emph{the stage-one threshold, \S\ref{s:margin}}
\item $\sqrt{H_\forall}\le\ell\le10\delta\sqrt{H_\forall}\log(2\delta H_\forall)$
 \hfill \textup{[S]} \cite[Cor.~2.2(ii), (C1)]{IUT4}, \emph{both sides}
\item $\dfrac{\ell+1}{24}N\log p>(\ell+5)\log(e^{*}_{\mathrm{mod}}\ell)$, both sides in nats
 \hfill \textup{[P]} from \cite[Thm.~1.10, Step (v)]{IUT4}
\item $e_{\mathrm b}=15\ell/\gcd(15\ell,\ord_p(q))\le p-2$
 \hfill \textup{[P]} $I^{*}=\varnothing$ in \cite[Prop.~1.4(iii)]{IUT4}
\item $\ell\nmid\ord_p(q)$ \hfill \textup{[S]} \cite[Def.~3.1(c)]{IUT1}
\end{enumerate}

\textup{[P]} \emph{Condition (4) is not a restriction on the source's estimate, and it
is worth being exact about this.}  \cite[Prop.~1.4(iii)]{IUT4} reads ``Let
$I^{*}\subseteq I$ be a subset such that for each $i\in I\setminus I^{*}$, it holds that
$p-2\ge e_i$'', so $I^{*}$ is \emph{chosen} and the inequality is required only outside
it; the bound then carries $4|I^{*}|/p$ and $\sum_{i\in I^{*}}\{3+\log(e_i)\}$.
\cite[Thm.~1.10, Step (v)]{IUT4} takes $I^{*}$ to be exactly the set of $i$ with
$e_{v_i}>p_{v_{\mathbb Q}}-2$, and by (R4) absorbs them into
$4(j+1)\,\ell^{*}_{\mathrm{mod}}$ --- the same term tabulated in \S\ref{s:slack}.  So
the estimate applies whether or not (4) holds; what (4) buys is the error-free form,
with $I^{*}=\varnothing$.

\textup{[P]} \emph{Two heights occur here and must be kept apart.}  Condition (C1) of
\cite[Cor.~2.2(ii)]{IUT4} is stated for $\log(q^{\forall})$, which keeps every
nonarchimedean place; write $H_{\forall}$ for it.  The height over the \emph{selected}
bad places, which drops those over $2$ and $\ell$, is a different number; write
$H_{\mathrm{sel}}$.  For $a=\pi^N$ over $\mathbb{Q}(\sqrt5)$, with $B=\Norm(1-a)$,
\[
 H_{\forall}=N\log p+\log B-\min(v_2(B),8)\log2,\qquad
 H_{\mathrm{sel}}(\ell)=N\log p+\log B-v_2(B)\log2-v_\ell(B)\log\ell .
\]
\textup{[P]} \emph{$H_{\mathrm{sel}}$ drops the place over that row's own $\ell$, not
over a fixed prime.}  The two heights agree only at $N=29$, where $\ell=7$ and
$v_7(B)=0$; at $N=114$, $\ell=37$ one has $v_{37}(B)=2$ and they differ by $2\log37$.
\emph{Condition (2) is tested against $H_{\forall}$}; both are printed below.

\begin{center}
\begin{tabular}{rrrrrcr}
\toprule
$N$ & $\ell$ & $H_{\forall}$ & $\sqrt{H_{\forall}}$ & $e_{\mathrm b}$ & (1)(2)(3)(4)(5) & $H_{\mathrm{sel}}$\\
\midrule
29  & 7  & $309.02$  & $17.58$ & 105 & \texttt{o X X o o} & $309.02$\\
\textbf{114} & \textbf{37} & $1217.45$ & $34.89$ & 185 & \texttt{\textbf{o o o o o}} & $1210.23$\\
160 & 53 & $1707.05$ & $41.32$ & 159 & \texttt{o o o o o} & $1704.28$\\
480 & 73 & $5132.24$ & $71.64$ & 73 & \texttt{o o o o o} & $5129.47$\\
480 & 79 & $5132.24$ & $71.64$ & 79 & \texttt{o o o o o} & $5129.47$\\
\bottomrule
\end{tabular}
\end{center}

\textup{[P]} \emph{Of the rows listed, $N=114$, $\ell=37$ is the first that passes all
five}, at the same prime $p=211$ and in the same family as \cite{First}.  \emph{The
conditions are not monotone in $N$ and do not stay satisfied}: in the same family
$N=116$, $\ell=37$ gives $e_{\mathrm b}=15\cdot37/\gcd(232,555)=555>209$ and fails (4).

\textup{[P]} \emph{Condition (5) has two parts and they are not on the same footing.}
At the designated place the order is $\ord_p(q)=2N$, and $\ell\nmid2N$ is an integer
check.  At the remaining bad places it asks that $v_{\mathfrak p}(1-a)$ be prime to
$\ell$, for which freedom of $\Norm(1-a)$ from $\ell$-th powers is \emph{sufficient but
not necessary} --- a valuation of $\ell+1$ is prime to $\ell$ and not $\ell$-th-power-free.  For the datum
\eqref{e:datum} of the present report that computation is shipped:
\path{scripts/imq_sieve7.cpp} runs the exhaustive sieve to $\lfloor C^{1/7}\rfloor$
(it is not the sieve for the datum of \cite{First}).  \emph{For the larger data of
\S\ref{s:region} and \S\ref{s:fourteen} it is not}: the cyclotomic factorization
$a-1=\prod_{d\mid N}\Phi_d(\pi)$ reduces the search, but on its own it excludes an
$\ell$-th power only for factors whose norm is small enough, and the certificate for the
reduced search is carried out by \path{scripts/powerfree_certificate.py}, which ships
with this report and runs in well under a second.  The rows above should be read
accordingly.

\begin{remark}[corrections]
\textbf{[P]} An earlier version of this table compared the two sides of condition (3)
\emph{in different units} --- the left in $\log p$, the right in nats --- which put
\texttt{X} in column (3) for $N=114$ and $N=160$ and made $N=480$ look like the first
point at which all five hold.  With both sides in nats --- equivalently, dividing both
by $\log p>0$, which is the form the code uses --- the first of the rows
listed at which all five hold is $N=114$ --- not every $N$ beyond it: $N=116$,
$\ell=37$ fails (4).  The
same earlier version used $e_{\mathrm b}=15\ell$ without the $\gcd$, omitted condition
(5), and reported a sieve analysis in a form the cyclotomic factorisation replaces by
a finite blockwise one.
All four are corrected here.  The rows $N=114$ and $N=160$ are due to an external
computation, which also supplied the certified values of $H$ used above for
calibration.
\end{remark}

\section{A curve carrying fourteen candidate levels}\label{s:fourteen}

\textup{[P]} The five conditions of \S\ref{s:region} do not act independently.  Raising
$N$ helps (1) and (3) but raises $H$, which by (2) forces $\ell$ up, which by (4)
forces $p$ up.  For a datum of the shape $a=\pi^N$ the tension resolves into a
\emph{window}: once (3), (4) and (5) hold, and once the upper half of (C1) is checked,
the admissible $\ell$ are the primes with
\[
  \sqrt{H_\forall}\ \le\ \ell\ <\ \tfrac{N}{3},
\]
the left end being the \emph{lower} half of the source's condition
\cite[Cor.~2.2(ii), (C1)]{IUT4} and the right end the threshold derived in
\S\ref{s:margin} --- \emph{the latter is ours, not the source's}.  \emph{Condition (C1)
is two-sided, and the window is not by itself equivalent to it}: we check the upper half
separately at each datum tabulated, and we do \emph{not} show that $\ell<N/3$ implies it
in general.  At the datum below that upper half,
$\ell\le10\delta\sqrt H\log(2\delta H)$ with
$\delta=2^{12}\cdot3^{3}\cdot5\cdot d$, where $d$ \emph{bounds} the degree in
$x_E\in U_X(\overline{\mathbb Q})^{\le d}$ \cite[Cor.~2.2]{IUT4}, is not binding: the
smallest admissible value here is $d=2=[\mathbb{Q}(a):\mathbb{Q}]$, giving
$\delta=1105920$ and a bound of about $2.2\times10^{10}$, and the bound only grows
with $d$; against a window whose right end is $157.33\ldots$.  \cite[Def.~3.1(c)]{IUT1} additionally requires
$\ell\ge5$ and $\ell$ prime to the residue characteristics of the bad places.  The
bad places lie above $p$ \emph{and} above every prime dividing $\Norm(1-a)$, so the
second condition is not settled by $\ell<p$; what settles it is that
$\ell\nmid\Norm(1-a)$, which we checked for each of the fourteen.  A datum is therefore interesting to the extent that this window
is wide.

\textup{[P]} Take
\[
  L=\mathbb{Q}(i),\quad \pi=41+26i,\quad \Norm(\pi)=2357,\quad a=\pi^{472},\quad
  E:y^2=x(x-1)(x-a),
\]
with $F=L(E[15])$ and $K_\ell=F(E[\ell])$.  Here $2357$ is prime and $\equiv1\pmod4$,
so it splits in $L$, and one place above it is multiplicative and one good.  Then
$\ord_p(q)=944$, the finite Tate height is $H=7324.7516\ldots$, and the window
$85.58\le\ell<157.33$ contains \emph{fourteen} primes:

\begin{center}
\begin{tabular}{rrrr}
\toprule
$\ell$ & $e_{\mathrm b}$ & $e_{\mathrm b}\le p-2$ & margin $s-3$ at $j=1$\\
\midrule
 89 & 1335 & yes & $205/89$\\
 97 & 1455 & yes & $181/97$\\
101 & 1515 & yes & $169/101$\\
103 & 1545 & yes & $163/103$\\
107 & 1605 & yes & $151/107$\\
109 & 1635 & yes & $145/109$\\
113 & 1695 & yes & $133/113$\\
127 & 1905 & yes & $91/127$\\
131 & 1965 & yes & $79/131$\\
137 & 2055 & yes & $61/137$\\
139 & 2085 & yes & $55/139$\\
149 & 2235 & yes & $25/149$\\
151 & 2265 & yes & $19/151$\\
\textbf{157} & \textbf{2355} & \textbf{yes} & $\mathbf{1/157}$\\
\bottomrule
\end{tabular}
\end{center}

\textup{[P]} At $\ell=157$ \emph{two upper constraints are simultaneously tight}:
$1/157$ is the last positive stage-one margin, since $472/3=157.33\ldots$; and
$e_{\mathrm b}=15\cdot157=2355=p-2$ exactly.  \emph{Both are upper constraints}; the
left wall of the window, $\sqrt{H_\forall}=85.58\ldots$, is what makes $89$ the first admissible
prime and has nothing to do with $157$.

\textup{[P]} \emph{What the fourteen are, exactly.}  They are the $\ell$ for which the
five numerical filters of \S\ref{s:region} hold at this curve.  \emph{They are not
fourteen verified initial $\Theta$-data}: no theorem here applies
\cite[Def.~3.1]{IUT1}(a)--(f) to each of them, the basis-free image argument of
\S\ref{s:arith} is carried out for the datum of \eqref{e:datum} and not for these, and
the order condition at the non-designated bad places is in the state described in
\S\ref{s:region}.  Two further remarks.  \emph{It is fourteen levels on one curve, not
fourteen curves}: $F_{\mathrm{mod}}$, $E$ and $p$ are fixed and only $\ell$ moves.
And the order condition of \cite[Def.~3.1(c)]{IUT1} at every bad place is what costs
the work here.  $\Norm(1-a)$ has $1592$ digits, so a direct $157$-th power bound runs
to about $10^{10.14}$.  The reduction uses \emph{nine blocks}: the cyclotomic factors
$\Phi_1,\Phi_2,\Phi_4,\Phi_8,\Phi_{59},\Phi_{118},\Phi_{236}$ kept whole, together with
\[
 G_{+}=\frac{X^{118}+i}{X^{2}-i},\qquad G_{-}=\frac{X^{118}-i}{X^{2}+i},
 \qquad G_{+}G_{-}=\Phi_{472},
\]
whose degrees $1,1,2,4,58,58,116,116,116$ sum to $472$.  \emph{It is $\Phi_{472}$ that
splits, not $\Phi_4$.}  Clearing the $\binom92=36$ pairwise gcds then reduces the
problem to a sieve small enough to run exhaustively, which is why the condition can
be settled for all fourteen levels at once rather than one at a time.  \emph{The digit
count, the factor count and the gcd count are checked by} \path{scripts/tables.py};
\emph{The sieve is shipped and runs}: \path{scripts/powerfree_certificate.py}
rebuilds the nine blocks from $\pi$ by exact Gaussian division --- each $\Phi_d(\pi)$
from $\prod_{e\mid d}(\pi^e-1)^{\mu(d/e)}$, with the product of all nine asserted equal
to $\pi^{472}-1$ --- then sieves each block norm exhaustively to its own $89$th-root
bound, the largest being $24{,}856$, and finds no $89$th power; $8{,}318$ primes are
tested in all.  It then checks the only way blockwise freeness could fail to give
freeness of $B$: the $36$ pairwise gcds show that just $2$ and $5$ meet two blocks, and
$v_2(B)=9$, $v_5(B)=2$, both below $89$.  Output in
\path{scripts/powerfree_certificate_output.json}; the same constants are recorded in
\path{lean/LargeImaginaryDatum.lean}, whose header states that having them does not by
itself prove power-freeness.  \emph{This establishes the sufficient condition, not
fourteen initial $\Theta$-data.}

\section{Coefficients printed in the source}\label{s:coeff}

\textup{[S]} Step (v) of the proof of \cite[Thm.~1.10]{IUT4} carries the coefficient
$j^2/(2\ell)$ on $\log(q)$ --- the \emph{unrooted} $q$, since
$\qbb=q^{1/(2\ell)}$ makes $\log(q^{j^2/(2\ell)})=j^2\log\qbb$, so on $\log\qbb$ the
coefficient would be $j^2$ --- and prints the average over
$j\in\{1,\dots,\ell^{*}\}$ first as $(2\ell^{*}+1)(\ell^{*}+1)/(12\ell)$ and then as
$(\ell+1)/24$.  At $\ell=7$:
\[
 \tfrac1{14},\ \tfrac27,\ \tfrac9{14}\quad\text{average}\ \tfrac13,
 \qquad\text{$q$-side}\ \tfrac1{2\ell}=\tfrac1{14},
 \qquad\text{ratio}\ \tfrac{12}{\ell(\ell+1)}=\tfrac3{14}.
\]
\textup{[P]} The two printed forms agree, and the average is $\ell(\ell+1)/12$ times
the stage-one coefficient.

\section{Local degrees and tensor orders}\label{s:orders}

\textup{[P]} This section is computed for the datum of \cite{First} at $211$; the
corresponding numbers for the datum of \eqref{e:datum} at $109$ differ, as noted in the
abstract.  At $211$ the good completion is unramified of degree $24$: the reduction
has $188$ points, trace $24$, and $X^2-24X+211$ has Frobenius of order $4,3,8$ modulo
$3,5,7$, with least common multiple $24$.  The bad completion has $e=105$, $f=2$.  So
$a_{\mathrm g}=24$, $a_{\mathrm b}=210$, and the tensor orders are
\[
 N_{(\mathrm g,\mathrm g)}=576,\qquad N_{(\mathrm g,\mathrm b)}=5040,
 \qquad N_{(\mathrm b,\mathrm b)}=44100 .
\]
\textup{[P]} \emph{The mixed tensor order is maximal and the all-bad one is not.}
Because the good side is unramified, $\mathcal O_{\mathrm g}\otimes\mathcal O_{\mathrm b}$
is the full ring of integers, with normalized gain $0$; for $(\mathrm b,\mathrm b)$ the
gain is $(e-1)/(2e)=52/105$ in units of $\log p$.  \textup{[S]}
\cite[Prop.~3.2]{IUT3} --- which is what \cite[Thm.~1.10]{IUT4} cites here, using
\cite[Thm.~3.11(i)(a)]{IUT3} only for the notation --- records that \emph{each of the $j+1$
factors, before tensoring}, is a direct sum indexed by the places of $\VVu$ over
$v_{\mathbb Q}$ --- $2$ of them here, $1$ over $\Fmod=\mathbb{Q}$.  \emph{This is not
the number of tuple components after tensoring} (that is $2^{j+1}$), nor the number of
field factors into which a local tensor product splits.

\section{Containers, seeds, and the floor slack}\label{s:container}

\textup{[S]} Step (v) works inside a container whose exponent is
$\lfloor\lambda\rfloor-|I|$.  \emph{The exponent alone does not name a set}: in
\cite[Prop.~1.4(iii)]{IUT4} it multiplies $\log_p(R_I^{\times})$, and the table below
records exponents, not lattices.  ``Seed'' in that table means the same thing rounded
the other way: $\lceil\lambda\rceil$ is the least integer exponent of the same lattice
$\log_p(R_I^{\times})$ that is at least $\lambda$, so the two columns are the upward
and downward discretisations of one real number against one reference lattice.  And
the bound of \cite[Prop.~1.4(iii)]{IUT4} takes a floor, leaving a slack.

\textup{[P]} \emph{The table is restricted to tuples whose final entry lies in
$\Vbad$}, so that $\lambda=j^2s$.  If the final entry is good then $\lambda=0$ by Step
(v) of \cite[Thm.~1.10]{IUT4}, and none of the numerical columns below applies.

\begin{center}
\begin{tabular}{ccccc}
\toprule
$j$ & $\lambda=j^2s\ (\underline v_j\in\Vbad)$ & exponent $\lfloor\lambda\rfloor-|I|$ & seed $\lceil\lambda\rceil$ & floor slack\\
\midrule
1 & $29/7$  & 2  & 5  & $6/7$\\
2 & $116/7$ & 13 & 17 & $3/7$\\
3 & $261/7$ & 33 & 38 & $5/7$\\
\bottomrule
\end{tabular}
\end{center}

\section{The slack in the averaged bound}\label{s:slack}

\textup{[S]} Step (v) of the proof of \cite[Thm.~1.10]{IUT4} carries, besides the
terms of \cite[Prop.~1.4(iii)]{IUT4}, a term $4\iota_{v_{\mathbb Q}}\ell^{*}_{\mathrm{mod}}$
with $\ell^{*}_{\mathrm{mod}}=\log(e^{*}_{\mathrm{mod}}\ell)$ and $\iota=1$ when
$p\le e^{*}_{\mathrm{mod}}\ell$; here
$e^{*}_{\mathrm{mod}}=2^{12}3^{3}5\,e_{\mathrm{mod}}$ with
$e_{\mathrm{mod}}\le d_{\mathrm{mod}}=2$.

\textup{[P]} For the datum of \cite{First}, $e^{*}_{\mathrm{mod}}\ell=7\,741\,440$, so
$\iota=1$ at $p=211$ and $\ell^{*}_{\mathrm{mod}}=2.964\log p$.  After the
procession-normalized average the term contributes $(\ell+5)\ell^{*}_{\mathrm{mod}}
=35.57\log p$, whereas the entire $q$-signal at $211$ is
$\tfrac{\ell+1}{24}\cdot29=9.667\log p$:
\[
 \boxed{\ \frac{(\ell+5)\,\ell^{*}_{\mathrm{mod}}}
   {\frac{\ell+1}{24}\cdot\frac{\ord_p(q)}{d_{\mathrm{mod}}}\cdot\log p}
 =3.6792605719\ldots\ }
\]
\textup{[P]} \emph{The $q$-term is written with the unrooted $q$}: here
$\ord_p(q)=58$ and $d_{\mathrm{mod}}=[\Fmod:\mathbb{Q}]=2$, so
$\tfrac{\ell+1}{24}\cdot\ord_p(q)/d_{\mathrm{mod}}=29/3$.  \emph{That this equals
$\tfrac{\ell+1}{24}\ord_p(\qbb)\,\ell$ is particular to $d_{\mathrm{mod}}=2$}, since
$\ord_p(\qbb)=\ord_p(q)/(2\ell)$.  \emph{Inside the averaged bound, the $\iota$-term
exceeds the $q$-term by a factor of nearly four.}  This compares two terms of one upper bound; it says nothing
about the distance from that bound to the quantity it bounds, and it is not a criticism
of the bound.  Note also that at $p=211$ one has $e_{\mathrm b}=105\le209$, so
$I^{*}=\varnothing$ is available here and the $\iota$-term is the price of the
\emph{uniform} form, not something forced by this datum.

\begin{remark}[correction]\label{r:box}
\textup{[P]} An earlier version printed this ratio without the $\log p$ in the
denominator, where the numerator is in nats and the denominator in units of $\log p$;
as printed it evaluates to $19.6908\ldots$, not $3.68$.  The intended quantity, with
both in nats, is the boxed one.
\end{remark}

\section{What the numbers suggest, with the reservations attached}\label{s:interp}

\textup{[P]} The tables above are meant to stand on their own, and each is checked by
\path{scripts/tables.py}.  We nonetheless set out what they appear to indicate, as
four readings, each paired with the reservation that limits it.  \emph{Nothing above depends on any of them}; where a
reading is echoed elsewhere --- (R1) in the abstract, (R2) in \S\ref{s:regime}, (R3) in
\S\ref{s:slack} --- it is as a remark, never as a step.

\begin{enumerate}[label=(R\arabic*),leftmargin=*]
\item \textbf{The admissible region is thin.}  For data of the shape $a=\pi^N$, and
for a candidate at which conditions (3), (4) and (5) of \S\ref{s:region} already hold,
the two remaining conditions reduce to $\sqrt{H_\forall}\le\ell<N/3$ \emph{together
with the upper half of} (C1) --- which is not binding at any datum tabulated here, and
which we have not shown to be non-binding in general (\S\ref{s:fourteen}).  So the room
available to numerical work on these estimates is set by \cite[Cor.~2.2(ii)]{IUT4} from
below and by the stage-one threshold from above.
\emph{Reservation:} the window is not by itself the admissible region --- $N=116$,
$\ell=37$ lies inside it and fails (4), with $e_{\mathrm b}=555>209$ --- and this is
computed for one shape of datum in two families.  A datum not of the form $a=\pi^N$
need not give a window of this kind, and we have not looked.

\item \textbf{The entry threshold is not, by itself, a criterion.}  Two curves in the
tabulated range clear it while requiring $I^{*}\ne\varnothing$ (\S\ref{s:regime}), so
``does the stage-one comparison say anything?'' is not settled by the threshold alone.
\emph{Reservation:} the sweep covers Cremona conductor below $10^4$ and a single
multiplicative place.  Larger conductors, and configurations with several bad places,
are untouched.

\item \textbf{At these parameters the $\iota$-term of the averaged bound dominates its
$q$-term}, by a factor of about $3.7$ (\S\ref{s:slack}).  \emph{Reservation:} this compares two terms \emph{inside} an
upper bound.  It is not a statement about the distance between that bound and the true
value --- which is what tightness would mean --- and it says nothing whatever about
whether the estimate holds.  A bound may be far from tight
and perfectly correct --- indeed \cite[Thm.~1.10]{IUT4} is stated for an arbitrary
datum, so slack at a particular one is what one should expect.

\item \textbf{These estimates can be evaluated on explicit input.}  Fourteen
candidate levels on one curve, together with the data of \cite{DH25} over $\mathbb{Q}$ and the family of
Proposition~\ref{p:family}, give enough explicit input for the local quantities to be
computed rather than discussed.  \emph{Reservation:} constructing a datum and
evaluating what the theory does with it are different tasks.  Ordinary normalised
log-volumes of the lattices involved, coefficients, and ratios of terms in the explicit
bounds are computed here; what is \emph{not} computed is the hull of the union of the
possible images of the actual $\Theta$-pilot, either side of the comparison of
\cite[Cor.~3.12]{IUT3}, or membership in the exceptional sets that the effective
statements exclude.
\end{enumerate}

\textup{[P]} One reading we deliberately do \emph{not} offer: that any of this bears on
whether \cite[Cor.~3.12]{IUT3} holds.  The quantities tabulated here are inputs to the
estimates, not their outputs, and no step of the present report examines the inference
those estimates rest on.

\begin{remark}[what this report draws on]\label{r:respect}
\textup{[P]} A report consisting of tables can read as an audit of the papers it draws
on.  It is not one, and the reason is worth stating precisely.  The conditions
in \S\ref{s:region} are of two kinds, and we keep them apart: (2) is
\cite[Cor.~2.2(ii)]{IUT4} and (4) is the condition under which
\cite[Prop.~1.4(iii)]{IUT4} may be applied with $I^{*}=\varnothing$; (5) is
\cite[Def.~3.1(c)]{IUT1}; and (1) and (3) are conditions imposed here.  The source's own conditions can be tabulated at
all only because they are stated precisely enough to be checked against a named curve,
and several errors in earlier drafts of this report were caught by those same
statements.  The explicit data of \cite{DH25} are what these questions are evaluated
against.  The programme described in \cite{LANA} is a separate formalization effort
that may prove complementary; we do not claim that its targets coincide with the
interface assumed here.  The exchange recorded in \cite{SS} and
\cite{Rpt} located the point at which the inference has to be understood; this report
does not enter it and takes no position on it.  The distinction that governs everything
above is between computing with a definition and understanding a proof: only the former
is attempted here.
\end{remark}

\section{Scope}\label{s:scope}

\noindent\textbf{Claimed.}  \eqref{e:datum} is an initial $\Theta$-datum with a
mixed local profile above $109$ and with $\Fmod$ imaginary quadratic; the family of
Proposition~\ref{p:family} exists.

\smallskip\noindent\textbf{Not claimed.}  That \cite[Thm.~A]{MFHMP} is affected or
that any estimate is; any conclusion about the $abc$ inequality; that any statement of
\cite{IUT1,IUT2,IUT3,IUT4} is false; that the hypothesis \SA{} of \cite{First} holds;
that the analysis of \cite{SS}, the responses of \cite{Rpt} or the findings of
\cite{LANA} bear on the datum.  Ordinary normalised log-volumes of the lattices
involved are computed here; what is not computed is the hull of the union of the
possible images of the actual $\Theta$-pilot, and what is not compared is the two
sides of \cite[Cor.~3.12]{IUT3}.  The relation of this datum to the estimates, if any, is left to a further
report; \emph{the construction stands whatever that relation turns out to be}.

\section{Reproduction}\label{s:repro}

\path{scripts/imq_checks.py} reproduces the exact arithmetic of \S\ref{s:arith}
(the six orbit comparisons, the exact $j$, the orders at $\pi$ and $\bar\pi$, the
norms and the Frobenius traces) using only the standard library.  \emph{It does not
reproduce the image argument}: the closure computation it also prints is an auxiliary
check on a chosen companion representative, labelled as such in its output, and the
basis-free argument is \path{lean/BasisFreeImage.lean}.  \path{scripts/imq_family.py} checks the three facts used in
Proposition~\ref{p:family} and the excluded density.
\path{scripts/tables.py} prints every table of
\S\ref{s:margin}--\S\ref{s:slack}, with the normalisations asserted rather than
printed, and exits zero; the fourteen levels of \S\ref{s:fourteen} and the sweep of
\S\ref{s:regime} are its Tables 11 and 12, each checking its own conclusion by
\texttt{assert} --- that the window holds exactly fourteen primes, that the two upper
constraints are simultaneously tight at $157$, and that at the two curves reaching the
threshold in the tabulated range the place must be put into $I^{*}$.  \emph{The counts
of the sweep itself are transcribed, not recomputed there.}  \path{scripts/imq_sieve7.cpp} performs the seventh-power-free sieve
\emph{for the cofactor $C$ of the datum of \eqref{e:datum} over $\mathbb{Q}(i)$}, not
for \cite{First}.  It is
self-contained --- the only operation it needs on $C$ is division by a machine word,
done on the decimal representation --- and runs in about a second.

\path{lean/BasisFreeImage.lean} carries the basis-free image argument of
\S\ref{s:arith}: twelve theorems, six of which depend on no axiom at all and six only
on \texttt{propext}, with no \texttt{sorryAx} and no \texttt{Classical.choice}.

Independently of those programs, two files under \path{lean/} ---
\path{ImaginaryQuadraticDatum.lean} and \path{ModSevenImage.lean} --- re-prove in the
Lean~4 kernel, importing nothing beyond \texttt{Init}, the integer identities and the
finite predicates of \S\ref{s:arith}: the norm factorisations, the valuations, the
point counts and traces $a_q$ at the witness places (recomputed there from the curve
rather than read off from this report), and the finite group computations.  \emph{The residue degrees are
not among them}; they are computed in \path{scripts/tables.py}, Table 13.

\textup{[P]} \emph{One doc-string in those files is out of date, deliberately.}  Seven
of the fourteen are pinned by SHA-256 to the bytes released with \cite{First}, and the
auditor treats any difference as an error, so a comment inside them cannot be edited
without giving up that guarantee.  \path{ModSevenImage.lean} accordingly still carries two descriptions we have since
retracted: its header presents the companion-matrix closure as an assertion about the
mod-$7$ image, and a doc-string says that the two places above a split $q$ ``give
different traces''.  Both are corrected here and neither enters a proof --- the file's
theorems are finite computations about a chosen representative, and the traces are
computed there from the curve.  The image argument is \path{BasisFreeImage.lean}, whose
own header states that it replaces the earlier one; the counterexample to the trace
claim, at $q=233$, is in \S\ref{s:arith}.
\emph{Four things are deliberately not in the kernel}, and the trust chain is only as
short as this list allows: the seventh-power-free sieve is not run there (what is
proved is that $34{,}869{,}670=\lfloor C^{1/7}\rfloor$, which is what makes the
external search exhaustive); the identification of these numeric predicates with the
actual Galois representation, places and cusps of the named curve is a field of the
interface below, not a theorem; the geometric arguments of \S\ref{s:geom} are not
verified by these programs; and \emph{Proposition~\ref{p:family} is formalised only in
part}.  \emph{That last one is worth stating plainly.}  Its proof combines finite
algebra with Hensel lifting, a sieve over all primes, a density bound and an
asymptotic; \emph{only the finite algebra is in the kernel}.
\path{lean/FamilyArithmetic.lean} carries three theorems --- that a prime dividing both
$r^2+1$ and $(r-1)^2+1$ divides $5$, its contrapositive, and the factorisation
$400=2^4\cdot5^2$ --- with no \texttt{sorryAx} and no \texttt{Classical.choice}.
\emph{The third is a numeric identity and not a discriminant theorem}: neither $f$ nor
the discriminant occurs in its statement, and $\operatorname{disc}(f)=400$ is computed
by ordinary means outside Lean.  \emph{The rest is prose.}  The density bound and the
infinitude of surviving $x$ are \emph{not formalised here} --- which is a statement
about this development, not a claim that they could not be.  That matters, because a gap in the prose part
passes every machine check in this bundle --- as one did: the root count needed
$q\nmid M$, the case $q=p$ was missing, and a later review, not the kernel, found it.
\emph{The theorem now in the kernel is exactly the one that was missing.}  Four further files treat \S\ref{s:geom} in the manner of an abstract interface: the cited results of anabelian geometry are the fields of a
structure \texttt{AnabelianInput}, each carrying its locator, and the corehood of the
datum's orbicurve is derived relative to any structure satisfying it.  Three things
there are proved outright, with no interface and no assumption: that $j(E_0)$ is
irrational, by an exact computation in which the unreduced numerator of the imaginary
part of $j(E_0)$ has $298$ digits (the reduced fraction reported in \S\ref{s:arith}
has $296$ and $234$) --- so that the
antecedent of \cite[Prop.~2.1]{MFHMP} is \emph{discharged} rather than assumed, the
four exceptional invariants being rational; that $\mathrm{SL}_2(\mathbb F_7)$ acts
transitively on the $24$ marked flags, which is what lets one global quotient and
cusp be matched at every bad place at once; and that $2\xi\ne\pm\xi$ for every
nonzero $\xi$, which separates the cusp $2\epsilon$ from $\epsilon$.  The file also
records two checks: a model in which the interface holds non-vacuously, and a model
showing that dropping \cite[Prop.~2.1]{MFHMP} makes the conclusion underivable.

Conditions (e) and (f) are derived in the same style.  The section
$\VV\to\mathbb V_{\mathrm{mod}}$ is proved \emph{bijective} rather than assumed, from
the transitivity above; the local conditions at the bad places and the arrow-marked
covers follow from seven further interface fields, each carrying its locator in
\cite{IUT1} or \cite{EtTh}.  Each of the seven comes with two checks: a model in
which its premise is actually inhabited, and a model in which that one field is
deleted, the other six still hold, and the conclusion fails.

\textup{[P]} \emph{Two kinds of arithmetic occur in these scripts and we do not call
both exact.}  The valuations, norms, residue degrees, point counts, ramification
indices, tuple coefficients and the right wall $N/3$ are integer or
\texttt{Fraction} computations and are exact.  The heights, the logarithms in
condition~(3) and the ratio of \S\ref{s:slack} are IEEE double evaluations of
$\log$ and $\sqrt{\ }$, and are exact only to display precision.  \emph{Where a
floating-point value decides a comparison, the margin is wide}: the left wall gives
$\sqrt{H_\forall}=85.58\ldots$ against the smallest admissible $\ell=89$, and
condition~(3) holds by a factor near $7.9$ across the window.  \emph{Those margins are
for the fourteen-level window only}; elsewhere they are smaller --- at $N=114$,
$\ell=37$ the gap to the left wall is $2.10\ldots$ and condition~(3) holds by a factor
of $1.31\ldots$ --- so we do not claim a single margin for the report.  No verdict here
is claimed to follow from an interval certificate, which we have not computed.

\path{scripts/lean_audit.py} checks \emph{every} declaration with
\texttt{\#print axioms}, including those the elaborator generates.  Its default scope
is all \emph{fourteen} files: the run of 2026-09-07 reports $263$ stated theorems and
$768$ declarations in all, of which $565$ depend on no axiom, $138$ on
\texttt{propext} and $65$ on \texttt{propext} and \texttt{Quot.sound}, and returns
\textsc{pass}.  Of those, \path{BasisFreeImage.lean}
contributes $12$ stated theorems and $37$ declarations, the twelve splitting $6$
axiom-free and $6$ on \texttt{propext}.  \emph{An earlier version of the auditor named
only eleven files and so did not cover that one}, which is the file carrying the
argument of \S\ref{s:arith}; the omission was found in a later review and is
corrected.  Seven of the fourteen --- the files of the first report --- are additionally
pinned by SHA-256, and carry $140$ stated theorems and $481$ declarations; their axiom
split $327/107/47$ is what the auditor prints under \texttt{-{}-through
GeometricModels}, which it labels a partial scope.  Nothing depends on
\texttt{sorryAx} or on \texttt{Classical.choice}; those two are not in the auditor's allowed set, so their
appearance is a failure rather than a note.  \emph{What the Lean files do not cover
is stated in their headers}: the cited anabelian results are assumed and not proved;
the links between the abstract interface and the concrete curve --- the $j$ of
\S\ref{s:arith}, and the finite predicates standing for places, flags and cusps ---
are named fields of the structures rather than theorems; and the sieve is not run,
what is proved being that $34869670=\lfloor C^{1/7}\rfloor$, which is what makes the
search exhaustive.  The geometric arguments of \S\ref{s:geom} are not
verified by these programs.


\begin{remark}[related work]
\textup{[S]} The LANA project reports that its plan is to work first toward the
formalization of the basic parts of anabelian geometry, and that its July 2026 interim
report provides a basis for eventual formalization in Lean rather than carrying one
out.  \textup{[P]} The Lean development described above does something different and
independent: it re-proves in the kernel the arithmetic of two explicit data against
\cite[Def.~3.1]{IUT1}, and states the results of anabelian geometry it uses as the
fields of an interface rather than proving them.  Were every field of that interface
proved, the conditional statements here would become unconditional.  We do not claim
that those fields are what LANA is formalizing: besides the anabelian theorems cited,
the interface also carries the identification of the abstract places, labels and cusps
with those of a named curve (\S\ref{s:repro}), which is a separate obligation.
\end{remark}

\section*{Disclosure}

Large language models were used substantially in locating passages, cross-checking
computations, drafting, and adversarial review; the author is responsible for the
claims.  The arithmetic of \S\ref{s:arith} and \S\ref{s:family} was verified by the
author's programs; the geometric arguments of \S\ref{s:geom} were drafted with
machine assistance against the sources cited there and then checked against those
sources.  The correction history of \cite{First} applies here as well, since the
present datum reuses its structure.  Several statements in earlier drafts of the
tabulations of \S\ref{s:margin}--\S\ref{s:slack} were withdrawn after an internal check: an incorrect ramification index, an omitted hypothesis of
\cite[Def.~3.1(c)]{IUT1}, a table applying a formula outside its domain, and a sieve
analysis reduced by a cyclotomic factorisation to a finite blockwise sieve.  The corrections are marked where they occur, except that the caveat on the four
$\mathbb{Q}$-rows of \S\ref{s:regime} stood only in the script until this version.

\begin{thebibliography}{MFHMP}
\bibitem[First]{First} T.~Urahigashi, \emph{An explicit initial $\Theta$-datum with
 mixed reduction at a split prime}, version 0.1.3-draft (2026),
 DOI \href{https://doi.org/10.5281/zenodo.22537342}{10.5281/zenodo.22537342}
 (all versions).
\bibitem[CanLift]{CanLift} S.~Mochizuki, \emph{The absolute anabelian geometry of
 canonical curves}, Doc. Math., Kazuya Kato's fiftieth birthday, Extra Vol.\
 (2003), 609--640.
\bibitem[EtTh]{EtTh} S.~Mochizuki, \emph{The \'etale theta function and its
 Frobenioid-theoretic manifestations}, Publ. Res. Inst. Math. Sci. \textbf{45}
 (2009), no.~1, 227--349.
\bibitem[IUT1]{IUT1} S.~Mochizuki, \emph{Inter-universal Teichm\"uller theory I},
 Publ. Res. Inst. Math. Sci. \textbf{57} (2021), 3--207.
\bibitem[IUT2]{IUT2} S.~Mochizuki, \emph{Inter-universal Teichm\"uller theory II},
 Publ. Res. Inst. Math. Sci. \textbf{57} (2021), 209--401.
\bibitem[IUT3]{IUT3} S.~Mochizuki, \emph{Inter-universal Teichm\"uller theory III},
 Publ. Res. Inst. Math. Sci. \textbf{57} (2021), 403--626.
\bibitem[IUT4]{IUT4} S.~Mochizuki, \emph{Inter-universal Teichm\"uller theory IV},
 Publ. Res. Inst. Math. Sci. \textbf{57} (2021), 627--723.
\bibitem[LANA]{LANA} Project LANA, \emph{Project LANA interim report on IUT
 theory}, July 2026, 50 pp.
\bibitem[Rpt]{Rpt} S.~Mochizuki, \emph{Report on discussions, held during the period
 March 15--20, 2018, concerning inter-universal Teichm\"uller theory}, RIMS, Kyoto
 University, February 2019 version, 45 pp.
\bibitem[SS]{SS} P.~Scholze, J.~Stix, \emph{Why $abc$ is still a conjecture},
 manuscript, July 16, 2018.
\bibitem[DH25]{DH25} T.~Dupuy, A.~Hilado, \emph{The statement of Mochizuki's
 Corollary 3.12: initial theta data}, arXiv:2004.13228v2 [math.NT], 19 June 2025,
 20~pp.  (Version~1, 28 April 2020, is a different and longer paper.)
\bibitem[DHPS]{DHPS} T.~Dupuy, A.~Hilado, \emph{Probabilistic Szpiro, baby Szpiro,
 and explicit Szpiro from Mochizuki's Corollary 3.12}, arXiv:2004.13108v2 [math.NT],
 29 April 2020.
\bibitem[MFHMP]{MFHMP} S.~Mochizuki, I.~Fesenko, Y.~Hoshi, A.~Minamide,
 W.~Porowski, \emph{Explicit estimates in inter-universal Teichm\"uller theory},
 Kodai Math. J. \textbf{45} (2022), no.~2, 175--236, DOI 10.2996/kmj45201.
\bibitem[Zhou]{Zhou} Z.-P.~Zhou,
\emph{The inter-universal Teichm\"uller theory and new Diophantine
results over the rational numbers.~I},
arXiv:2503.14510v1 [math.NT], submitted 8 March 2025, 41~pp.
\bibitem[Cremona-data]{CremonaData} J.~E.~Cremona,
\emph{Elliptic curve data}, \texttt{ecdata},
file \texttt{allcurves.00000-09999},
\url{https://github.com/JohnCremona/ecdata}, commit
\href{https://github.com/JohnCremona/ecdata/tree/052d21ef383c1416dd98017c9ae55158822d29ac}{\texttt{052d21ef383c}}
(14 July 2014).
\end{thebibliography}

\end{document}
